\section{Classical computation foundations}

Quantum circuits build on ideas from classical digital logic. Understanding
bits, Boolean functions, truth tables, and reversible computation makes the
QPU gate model much easier to use.

\subsection{Classical bits and logical state}

A classical bit has one definite logical value:
\[
b\in\{0,1\}.
\]
Physical computers may represent that value with voltage, charge, magnetic
orientation, or another measurable quantity. Digital logic ignores most
physical detail and works with the abstraction 0 or 1.

A wire carries a bit from one operation to another. A \emph{gate} computes a
function of one or more incoming bits. For example, NOT has one input and one
output:
\[
\operatorname{NOT}(0)=1,\qquad \operatorname{NOT}(1)=0.
\]
In QPU source, a token names the logical wire. On computational-basis states,
\code{X -I Q -O Q} has the same measured input/output relationship as
classical NOT.

\subsection{Boolean algebra}

Boolean algebra describes functions whose values are 0 and 1. This guide uses:
\begin{center}
\begin{tabularx}{\textwidth}{>{\bfseries}l c l X}
\toprule
Operation & Symbol & Spoken form & Meaning\\
\midrule
NOT & $\neg A$ & not A & 1 when A is 0; 0 when A is 1.\\
AND & $A\land B$ & A and B & 1 only when both inputs are 1.\\
OR & $A\lor B$ & A or B & 1 when either or both inputs are 1.\\
XOR & $A\oplus B$ & A exclusive-or B & 1 when exactly one input is 1.\\
NAND & $\neg(A\land B)$ & not-and & The complement of AND.\\
\bottomrule
\end{tabularx}
\end{center}

For bits, XOR is addition modulo two:
\[
0\oplus0=0,\quad 0\oplus1=1,\quad
1\oplus0=1,\quad 1\oplus1=0.
\]
There is no carry in modulo-two addition. Ordinary binary addition handles the
carry separately, which is why a half-adder produces both Sum and Carry.

\subsection{Truth tables as complete finite specifications}

A truth table lists a function's output for every possible input pattern. With
$n$ binary inputs there are $2^n$ patterns. AND, OR, XOR, and NAND can be
compared in one table:
\begin{center}
\begin{tabular}{cc|cccc}
\toprule
$A$&$B$&$A\land B$&$A\lor B$&$A\oplus B$&$\neg(A\land B)$\\
\midrule
0&0&0&0&0&1\\
0&1&0&1&1&1\\
1&0&0&1&1&1\\
1&1&1&1&0&0\\
\bottomrule
\end{tabular}
\end{center}
For a deterministic combinational circuit, this table completely specifies
the classical input/output behavior. It does not reveal the circuit's internal
gate arrangement, propagation time, power consumption, or quantum phase.

\subsection{Combinational logic and stored state}

A \emph{combinational} output depends only on current inputs. Adders and the
Boolean gates above are combinational. A \emph{sequential} circuit also depends
on stored state from an earlier step, as in counters and latches.

QPU cycle suffixes can document successive stages, and INCREASECYCLE marks a
workspace boundary, but the compiled circuit is still one ordered gate
sequence. \code{SAVE_STATE} and \code{LOAD_STATE} copy and restore the simulator
state vector. They are not quantum gates. When modeling stateful behavior
inside the unitary circuit, expose previous state as PARAMS and return next
state through RETURNVALS.

\subsection{Why ordinary classical gates can lose information}

A function is reversible only if each output identifies exactly one input.
Classical NOT is reversible: seeing output 1 proves the input was 0. AND is not:
output 0 could have come from 00, 01, or 10.

\[
(A,B)\xrightarrow{\mathrm{AND}}A\land B
\]
maps four input patterns onto only two output values, discarding information.
Quantum evolution through closed-system gates must be unitary and therefore
reversible. A quantum/reversible AND preserves both inputs and mutates another
target:
\[
(A,B,t)\longmapsto(A,B,t\oplus(A\land B)).
\]
Knowing the output triple lets the transformation be run backward to recover
the input triple. This one-to-one property is one of four conditions a
quantum gate must satisfy to be reversible; \emph{Requirements for a
reversible gate}, in the Theory Guide's mathematical foundations chapter,
states all four together.

\begin{beginnerbox}[How to read $t' = t\oplus(A\land B)$]
\begin{itemize}
  \item $A$ and $B$ are \textbf{controls}. They are only read. After the gate
  they still hold exactly what they held before, which is why the input can
  always be recovered.
  \item $t$ is the \textbf{target}, the output workspace. It is the only wire
  that changes. The prime in $t'$ means ``$t$ after the gate''.
  \item $\oplus$ means ``flip if''. The target flips when $A\land B$ is 1 and
  stays put otherwise.
  \item Start the target at 0 and it ends holding $A\land B$, the ordinary AND
  answer. Start it at 1 and it ends holding NAND instead.
\end{itemize}
In the Circuit builder, $A$ is \textbf{Control A}, $B$ is
\textbf{Control B}, and $t$ is \textbf{Target particle}.
\end{beginnerbox}

\subsection{Controls, targets, and workspace}

These terms connect classical logic to QPU syntax:
\begin{description}
  \item[Control] A wire whose value decides whether a target operation occurs.
  In \code{CNOT -I A -O T}, A is the control.
  \item[Target] The wire changed by an operation. T is the target in that CNOT.
  \item[Workspace or ancilla] An extra wire introduced to hold an intermediate
  or output while preserving inputs. It is commonly initialized to $\ket0$.
  \item[Garbage] Intermediate information left in workspace after the desired
  answer has been computed. Reversible algorithms often \emph{uncompute}
  temporary gates to return garbage wires to zero.
  \item[Uncomputation] Applying inverse operations in reverse order to clean
  workspace without erasing information irreversibly.
\end{description}

The QPU language calls logical wire names tokens. \code{CREATETOKEN} reserves
workspace, SET prepares it, \code{-I} lists inputs/controls, and \code{-O}
names targets.

\subsection{Fan-out, copying, and the no-cloning boundary}

Classical logic freely copies a bit onto several wires. If the source is a
known basis bit and targets begin at zero, CNOT reproduces that basis value:
\[
\ket{x}\ket0\xrightarrow{\mathrm{CNOT}}\ket{x}\ket{x},
\qquad x\in\{0,1\}.
\]
For a superposition $\alpha\ket0+\beta\ket1$, the same circuit produces
\[
\alpha\ket{00}+\beta\ket{11},
\]
not
\[
(\alpha\ket0+\beta\ket1)\otimes
(\alpha\ket0+\beta\ket1).
\]
The first state is generally entangled; the second would be two independent
copies and contains cross terms. The no-cloning theorem forbids a universal
gate that copies every unknown quantum state. QPU aliases similarly give two
names to one wire; they do not create a physical copy.

\subsection{Deterministic, probabilistic, and quantum models}

\begin{tabularx}{\textwidth}{>{\bfseries}p{0.19\textwidth}XXX}
\toprule
Model & State description & Transformation & Observation\\
\midrule
Deterministic classical & One definite bit string & Boolean function &
The definite output bits\\
Probabilistic classical & Probabilities over bit strings & Stochastic update &
A sampled string\\
Pure-state quantum & Complex amplitudes over basis strings & Unitary matrix &
A sampled string after magnitude-squared probabilities and collapse\\
\bottomrule
\end{tabularx}

Negative and imaginary amplitudes are not negative or imaginary probabilities.
They carry phase. Two computational paths can add constructively or cancel
destructively before probabilities are calculated. This interference has no
counterpart in an ordinary classical random-choice table.

\subsection{Classical and QPU operation map}

\begin{center}
\begin{tabularx}{\textwidth}{>{\bfseries}l l X}
\toprule
Classical idea & QPU operation & Important difference\\
\midrule
NOT & X or NOT & Acts linearly on superpositions; X is unitary.\\
XOR into target & CNOT & Preserves the control and can entangle.\\
AND into target & CCNOT or AND & Preserves controls; target is XOR-mutated.\\
Swap variables & SWAP & Exchanges complete quantum wire states.\\
Random bit & H then MEASURE & H first creates coherent amplitudes, not an
ordinary hidden random choice.\\
Read a bit & MEASURE & Sampling collapses a superposition and can disturb
entanglement.\\
Initialize zero & SET token 0p & Compiles to state preparation/reset behavior.\\
Function output & RETURNVALS & Declares an interface; does not itself copy or
measure a wire.\\
\bottomrule
\end{tabularx}
\end{center}

